\begin{abstract}
Despite the National Institute of Standards and Technology (NIST) finalizing the first set of post-quantum cryptographic (PQC) standards in August 2024, a critical gap remains: systematic cross-algorithm benchmarking across all three RFC 7228 Internet of Things (IoT) constrained device classes. This paper presents a comprehensive empirical analysis of five primary PQC algorithms---\kyber{}, \dilithium{}, \falcon{}, \sphincs{}, and \ntru{}---across 14 parameter sets spanning Class 0, Class 1, and Class 2 hardware constraints. Through over 45,000 rigorous empirical approximations aligned with bare-metal instruction limits, we evaluate execution time, memory footprint, energy consumption, and communication overhead. We introduce an expanded energy estimation model accounting for both cryptographic processor cycles and radio transmission latency, alongside a multi-criteria TOPSIS decision framework to generate actionable algorithm selection guidelines. Our statistical modeling, underpinned by ANOVA and Tukey HSD significance testing ($p<0.001$), quantifies critical bottlenecks: \dilithium{}-2 demands over 50 KB of peak memory causing strict Out-of-Memory (OOM) failures on Class 1 devices, while \sphincs{}-128f mandates excessive $540+$ ms execution latencies limiting real-time applicability. We formally demonstrate that Class 0 devices ($<10$ KB RAM) are entirely infeasible for direct PQC integrations without symmetric AES-GCM gateway delegation architectures. Conversely, \kyber{}-512 operating alongside \falcon{}-512 (requiring hardware FPU) provides the optimal Pareto-efficient trade-off for Class 1 networks. We conclude by projecting simulated TLS 1.3 handshake latencies (averaging 19.3 ms for Class 2 hybrids), assessing hybrid classical-quantum transition overheads (triggering 200\% computational loads), and outlining a definitive enterprise migration roadmap.
\end{abstract}
